\chapter{Введение: память как единый субстрат}

\section{Зачем вообще граф}

У большинства современных диалоговых агентов память устроена плоско: история
сообщений, иногда векторная база, куда складываются «воспоминания» как
независимые записи, и retrieval по косинусной близости к текущему запросу. Это
работает для ассистента, который обслуживает разрозненные запросы, но плохо
описывает существо, у которого есть биография, характер, меняющееся настроение
и убеждения, способные оказаться ошибочными. У человека память не список; она
сеть. Одно воспоминание тянет за собой другое; факт связан с тем, откуда он
узнан; недавняя новость медленно перекрашивает старые выводы; забытое не
исчезает мгновенно, а тускнеет и иногда всплывает.

Поэтому центральная ставка проекта Рин — \emph{память как граф}. Но сам по себе
выбор структуры данных ничего не решает. Интересна не «база в виде графа», а то,
какие процессы по этому графу текут и каким законам они подчиняются. Эта
монография — про процессы.

\section{Главный тезис}

\begin{intuition}[Одна фраза, которую стоит держать в голове до конца]
Память Рин — это \textbf{один типизированный темпоральный граф с метаданными на
узлах и рёбрах}, а всё остальное — \emph{операторы над этим субстратом}:
retrieval, пропагация уверенности, забывание и консолидация, ночные обходы,
настроение, слой личности. Если довести до чистоты \emph{схему субстрата}, то
все прочие механизмы становятся модулями, читающими и пишущими одни и те же
поля. Если схему не довести — каждый модуль потащит свою полусовместимую модель
памяти, и они не срастутся.
\end{intuition}

Эта установка определяет и порядок изложения, и порядок разработки. Сначала мы
описываем субстрат (глава~\ref{ch:substrate}) и его временную многослойность
(глава~\ref{ch:timescales}). Затем — динамические процессы: настроение
(глава~\ref{ch:mood}), обход графа (глава~\ref{ch:walker}), работу с
уверенностью и ревизией убеждений (главы~\ref{ch:epistemics}
и~\ref{ch:propagation}). Отдельная глава~\ref{ch:topology} посвящена тому, как
\emph{топология} самого графа незаметно управляет всеми этими потоками. Затем мы
переходим к слою личности и голоса (главы~\ref{ch:personality}
и~\ref{ch:voice}), к вопросу о том, в каком смысле всё это является
\emph{обучением без обучения} (глава~\ref{ch:learning}), и, наконец, к программе
экспериментов на суррогатном графе (глава~\ref{ch:experiment}).

\section{Четыре кита и почему их нельзя тюнить порознь}

Полезно заранее назвать четыре процесса, вокруг которых всё вращается:

\begin{enumerate}
\item \textbf{Правильный retrieval} — как из графа достаётся релевантное, с
учётом запроса и текущего состояния.
\item \textbf{Правильная пропагация уверенности} — как новый факт меняет
уверенность соседних узлов, не устраивая ни ступора, ни лавины.
\item \textbf{Правильное забывание и консолидация} — что тускнеет, что
укрепляется, что сливается в обобщения.
\item \textbf{Правильные ночные обходы} — фоновая реорганизация и починка,
которую агент делает «во сне».
\end{enumerate}

Важнейшее наблюдение: эти четыре кита \emph{не независимы}. Они читают и пишут
одни и те же поля активации и уверенности на одном и том же графе. Забывание
меняет топологию; изменённая топология меняет retrieval; изменённый retrieval
меняет то, что пропагируется; пропагация меняет то, что найдёт ночной обходчик.
Поэтому их нельзя настраивать по отдельности — их надо со-настраивать, и именно
поэтому (глава~\ref{ch:experiment}) нужен управляемый стенд с заранее известной
истиной: без него взаимодействие операторов просто не измерить.

\section{Родство с математикой DTI: одна спина, два применения}

Эта монография выросла рядом с другим исследованием автора — магистерской
работой о предсказании аффинности связывания белок–лиганд средствами
графового и геометрического глубокого обучения. На первый взгляд это далёкие
темы. На деле у них общая математическая спина.

Консолидация памяти (переход от подробных эпизодов к сжатым смыслам) — это
\emph{огрубление графа} (\textit{coarse-to-local}): какой сигнал сохраняется при
переходе от мелкого масштаба к крупному, а какой теряется. Это в точности
центральный вопрос DTI-работы, только там он про interface-релевантный сигнал
взаимодействия, а здесь — про то, что выживает при консолидации воспоминаний.
Дальше по тексту мы увидим, что over-squashing (удушение дальней связи узким
местом графа), спектральные представления, персистентный лапласиан и
эффективное сопротивление — инструменты, которые работают в обе стороны.
Осознание этого родства — не поэтическая вольность: те же операторы, только
приложенные к разным графам. Мы будем возвращаться к этой связке, когда она
даёт конкретную пользу.

\section{Для кого этот текст}

У монографии два адресата, и оба важны.

Первый — \emph{автор через две недели}, вернувшийся после напряжённой работы над
другими задачами (магистерская, беттинг) и забывший половину деталей. Для него
текст должен восстанавливать замысел \emph{с нуля}, а не намекать на него.

Второй — \emph{коллеги по команде}, которые в проектировании этой памяти пока не
участвовали и «не в теме». Для них текст должен вводить каждое понятие явно, до
первого употребления, и объяснять не только «как», но и «зачем».

Отсюда жёсткое правило изложения: \textbf{ни один объект или механизм не
используется раньше, чем определён}. Глава~\ref{ch:architecture} даёт систему
целиком одним взглядом; глава~\ref{ch:definitions} — словарь всех основных
объектов; и лишь затем идут детальные разборы. Примеры сознательно берутся из
жизни Рин (стрим, зрители, двое пользователей проекта, бытовые ситуации), чтобы
механизм можно было «пощупать» на конкретной гипотетической сцене.

\section{Как читать монографию}

Изложение намеренно избыточно. Там, где метод опирается на инженерную или
научную идею со стороны (труд-система поддержки истинности, самоорганизованная
критичность, теория оценивания эмоций, ACT-R и прочее), мы кратко описываем и
саму идею, чтобы читатель не был вынужден держать в голове внешний контекст.
Формулы служат для точности, а врезки «Интуиция» и «Ловушка» — для смысла и для
предупреждения о типичных ошибках. Ссылки на литературу собраны в
главе~\ref{ch:references}; часть из них — работы из библиографий рабочих волтов
проекта (RIN и RIN–DTI), это отмечено.
